\documentclass[a4paper,11pt]{article}

\usepackage[utf8]{inputenc}
\usepackage{amsmath,amssymb}
\usepackage[a4paper,margin=2.5cm]{geometry}
\usepackage{enumitem}

\title{Experiment Day:\\Harmonic Oscillators}
\author{}
\date{}

\begin{document}

\maketitle

\section{Starting point: spring and simple pendulum}

install spring and simple pendulum \\ 
\textbf{Discussion}
\begin{itemize}
    \item name similarities and differences of both systems
    \item What is the equilibrium position?
    \item What is the amplitude?
    \item What is one complete oscillation?
    \item Which system oscillates faster?
    \item What could affect the frequency?
\end{itemize} 
\\  
 \textbf{Predictions - important quantities for frequency }
\begin{itemize}
    \item spring pendulum: mass $m$, spring constant $k$, amplitude $A$
    \item simple pendulum: mass $m$, string length $l$, amplitude $A$
\end{itemize} 

\paragraph{Good Scientific Practice}

\begin{itemize}
    \item Change only \textbf{one parameter at a time}.
    \item Repeat every measurement several times.
    \item minimize error (e.g. measure 10 periods)
    \item Record your data immediately.
    \item Use the same measuring method throughout the experiment.
    \item Do not ignore measurements that do not fit your expectations. Instead, think about possible sources of error.
\end{itemize}
\section{General Procedure}

In science, experiments are performed to answer a specific question. To obtain reliable results, it is important to work carefully and systematically.

For every experiment, follow these steps:

\begin{enumerate}
    \item \textbf{State the research question.}\\
    What do you want to investigate?

    \item \textbf{Make a prediction (hypothesis).}\\
    What do you expect to happen? Explain your reasoning.

    \item \textbf{Identify the variables.}\\
    Decide which parameter you want to change (\emph{independent variable}) and which quantity you will measure (\emph{dependent variable}). Keep all other parameters constant to ensure a fair comparison.

    \item \textbf{Perform the measurements.}\\
    Measure the time for 10 oscillations instead of only one to reduce the influence of your reaction time. Repeat each measurement at least three times and calculate the average value.

    \item \textbf{Calculate the period.}
    \[
    T=\frac{t_{10}}{10}
    \]

    \item \textbf{Calculate the frequency.}
    \[
    f=\frac{1}{T}
    \]

    \item \textbf{Analyze your results.}\\
    Compare the measurements for the different parameter values. Look for trends and relationships.

    \item \textbf{Draw a conclusion.}\\
    Was your hypothesis correct? Which parameter influences the oscillation? Can you explain your observations?
\end{enumerate}
\section{Experiments}

\subsection*{Experiment 1 -- Spring Oscillator: Influence of the Mass}

\textbf{Question}
\begin{itemize}
\item How does the mass affect the oscillation?
\end{itemize}

\textbf{Materials}
\begin{itemize}
\item spring
\item stand with clamp
\item different masses
\item stopwatch 
\end{itemize}

\textbf{Procedure}
\begin{itemize}
\item Keep the spring constant.
\item Keep the amplitude approximately constant.
\item Change only the mass.
\item Measure the time for 10 oscillations.
\end{itemize}

\subsection*{Experiment 2 -- Spring Oscillator: Influence of the Spring}

\textbf{Question}
\begin{itemize}
\item How does the spring stiffness affect the oscillation?
\end{itemize}

\textbf{Materials}
\begin{itemize}
\item different springs
\item one mass
\item stand
\item stopwatch
\end{itemize}

\textbf{Procedure}
\begin{itemize}
\item Keep the mass constant.
\item Change only the spring.
\item Measure the time for 10 oscillations.
\end{itemize}

\subsection*{Experiment 3 -- Spring Oscillator: Influence of the Amplitude}

\textbf{Question}
\begin{itemize}
\item Does the amplitude affect the oscillation?
\end{itemize}

\textbf{Materials}
\begin{itemize}
\item spring
\item one mass
\item stopwatch
\end{itemize}

\textbf{Procedure}
\begin{itemize}
\item Keep the spring and the mass constant.
\item Change only the amplitude.
\item Measure the time for 10 oscillations.
\end{itemize}

\subsection*{Experiment 4 -- Simple Pendulum: Influence of the Length}

\textbf{Question}
\begin{itemize}
\item How does the pendulum length affect the oscillation?
\end{itemize}

\textbf{Materials}
\begin{itemize}
\item pendulum bob
\item string
\item stand
\item ruler
\item stopwatch
\end{itemize}

\textbf{Procedure}
\begin{itemize}
\item Keep the mass constant.
\item Change only the string length.
\item Use small amplitudes.
\item Measure the time for 10 oscillations.
\end{itemize}

\subsection*{Experiment 5 -- Simple Pendulum: Influence of the Mass}

\textbf{Question}
\begin{itemize}
\item Does the mass affect the oscillation?
\end{itemize}

\textbf{Materials}
\begin{itemize}
\item strings of equal length
\item different masses
\item stopwatch
\end{itemize}

\textbf{Procedure}
\begin{itemize}
\item Keep the string length constant.
\item Keep the amplitude constant.
\item Change only the mass.
\item Measure the time for 10 oscillations.
\end{itemize}

\subsection*{Experiment 6 -- Simple Pendulum: Influence of the Amplitude}

\textbf{Question}
\begin{itemize}
\item Does the amplitude affect the oscillation?
\end{itemize}

\textbf{Materials}
\begin{itemize}
\item pendulum
\item stopwatch
\end{itemize}

\textbf{Procedure}
\begin{itemize}
\item Keep the length and mass constant.
\item Change only the amplitude.
\item Measure the time for 10 oscillations.
\end{itemize}

\section{Discussion}

Discuss the following questions:

\begin{itemize}
    \item Which parameters influence the frequency?
    \item Which parameters have no influence?
    \item Which results surprised you?
    \item How are the spring oscillator and the simple pendulum similar?
    \item How do they differ?
\end{itemize}


\section{Similarities and Differences}

\subsection*{Similarities}

Both the spring oscillator and the simple pendulum share many common properties:

\begin{itemize}
    \item Both are \textbf{harmonic oscillators} (for small oscillations).
    \item Both oscillate around an \textbf{equilibrium position}.
    \item Both are characterized by an \textbf{amplitude}, a \textbf{period}, and a \textbf{frequency}.
    \item Both can be described using \textbf{sine and cosine functions}.
    \item In both systems, the \textbf{restoring force} always points back towards the equilibrium position.
    \item Because of \textbf{inertia}, the object overshoots the equilibrium position, leading to continuous oscillations.
\end{itemize}

\subsection*{Differences}

Although both systems behave similarly, the restoring force and the parameters determining the frequency are different.

\paragraph{Restoring force}

\begin{itemize}
    \item \textbf{Spring oscillator:} The restoring force is provided by the spring and is described by Hooke's law,
    $
    F=-kx.
    $

    \item \textbf{Simple pendulum:} The restoring force is provided by gravity.
\end{itemize}
 \\ 

\paragraph{Frequency} \\ 

\\ 
The frequency of a spring oscillator is $f=\frac{1}{2\pi}\sqrt{\frac{k}{m}},$ where $k$ is the spring constant and $m$ is the mass.

The frequency of a simple pendulum is$f=\frac{1}{2\pi}\sqrt{\frac{g}{l}},$ where $g$ is the gravitational acceleration and $l$ is the pendulum length.

\paragraph{Comparison}

\begin{itemize}
    \item For the \textbf{spring oscillator}, the frequency depends on the \textbf{mass} and the \textbf{spring constant}.
    \item For the \textbf{simple pendulum}, the frequency depends on the \textbf{pendulum length} and the \textbf{gravitational acceleration}.
    \item Increasing the mass decreases the frequency of a spring oscillator, whereas the mass has \textbf{no influence} on the frequency of a simple pendulum.
    \item A stiffer spring increases the frequency, while a longer pendulum decreases the frequency.
\end{itemize}

\paragraph{Key Takeaway}

Although the spring oscillator and the simple pendulum look quite different, they follow the same mathematical model of harmonic motion. The main difference is which physical quantities determine their frequency.

\end{document}
